\section{Functional Requirements}

Functional requirements define the explicit services, interactions, and behaviours that the system must provide in order to satisfy stakeholder expectations. In the proposed AI-Driven Lab Exam Monitoring and Management System, these requirements are interconnected because a practical lab exam must move in a controlled sequence from configuration to execution, monitoring, submission, and reporting.

Another important characteristic of these requirements is that they reflect academic control as well as software capability. The purpose of the system is not only to automate routine work, but also to provide exam governance in a form that is transparent, auditable, and fair. The following requirements are therefore written in operational sequence.

\begin{table}[H]
\centering
\caption{Functional Requirements}
\label{tab:functionalrequirements}
\renewcommand{\arraystretch}{1.2}
\begin{tabular}{lp{10.6cm}}
\hline\hline
\textbf{Requirement ID} & \textbf{Description} \\
\hline\hline
FR1 & Admin can manage user roles and system accounts. \\
FR2 & User can log in to the system. \\
FR3 & System shall bind the user workstation via Hardware ID (HWID). \\
FR4 & Teacher can create and configure exam templates and questions. \\
FR5 & Teacher can define the allowed application list for the exam. \\
FR6 & Teacher can allocate seats and manage student eligibility. \\
FR7 & Student can run system diagnostics and mock tests before the exam. \\
FR8 & Student can view scheduled exams and instructions on the dashboard. \\
FR9 & System shall lock down the student workstation to restrict external applications. \\
FR10 & Student can start and conduct the exam inside the secure environment. \\
FR11 & System shall stream real-time screen captures and focus telemetry. \\
FR12 & Teacher can monitor student activities and focus timelines live. \\
FR13 & Student can submit exam answers, and system shall perform auto-save backups. \\
FR14 & System shall compile and grade code submissions using test cases. \\
FR15 & System shall grade theory answers using AI semantic evaluation. \\
FR16 & System shall check code submissions for plagiarism and similarity. \\
FR17 & Admin and Teacher can view transaction audit logs. \\
FR18 & Teacher can generate and download analytical PDF exam reports. \\
\hline\hline
\end{tabular}
\end{table}


\section{Non-Functional Requirements}

Non-functional requirements define the quality conditions under which the functional services must operate \cite{iso25010}. In a project concerned with secure and fair academic assessment, non-functional properties are not secondary. If the system is difficult to use, invigilators may ignore it during the exam. If it performs poorly, the dashboard may become operationally irrelevant. If it is insecure, confidence in the entire assessment process may decline. For these reasons, the proposed system must satisfy multiple quality dimensions simultaneously.

\subsection{Software Quality Attributes}

The main software quality attributes for the proposed system are security, reliability, usability, performance, maintainability, and data integrity. Security is essential because the platform manages identities, exam schedules, student records, and potentially sensitive monitoring evidence. Reliability is equally important because a lab exam cannot pause simply because the system behaves inconsistently. Usability matters because administrators and teachers must interact with the platform under time pressure, while students need a clear and low-confusion exam interface.

Performance is also important because the system must respond quickly during login, attendance verification, exam launch, live monitoring, and file submission. Maintainability is relevant because the project may later be refined, extended, or integrated with broader institutional workflows. Data integrity remains critical throughout because schedules, attendance status, incident logs, results, and reports must remain consistent and auditable.

\subsection{Other Non-Functional Requirements}

Other relevant non-functional requirements include portability, privacy sensitivity, compatibility with university lab settings, and support for structured backup practices. The proposed system should operate within the technological realities of academic computer labs, where hardware and operating environments may vary. It should also preserve logs and exam records in a consistent manner so that important information is not lost due to technical disruption or later review needs.

Privacy considerations are particularly important when user activity, event logs, or submission-related evidence are retained. The system should therefore support responsible access control and discourage unnecessary exposure of sensitive student-related information. In addition, the platform should remain easy to document and explain within an academic environment, because FYP success depends not only on implementation but also on clarity of presentation and justification.

The non-functional requirements show that the platform must be technically sound and institutionally dependable at the same time. Quality in this project cannot be judged only by whether a screen works. It must also be judged by whether the platform remains trustworthy when used under real examination pressure.

\begin{table}[H]
\centering
\caption{Non-Functional Requirements}
\label{tab:nonfunctionalrequirements}
\begin{tabular}{lp{10cm}}
\hline\hline
\textbf{Category} & \textbf{Requirement} \\
\hline\hline
Security & The system shall secure all accounts and data. \\
Reliability & The system shall ensure data reliability. \\
Usability & The system shall be easy to use for everyone. \\
Performance & The system shall respond quickly to all actions. \\
Maintainability & The system shall be easy to update and maintain. \\
Privacy & The system shall protect all sensitive information. \\
Compatibility & The system shall work in university lab settings. \\
Backup and Recovery & The system shall backup and recover all data. \\
\hline\hline
\end{tabular}
\end{table}

\section{Requirement Gathering Techniques Used}

The requirements for the proposed system were gathered through a multi-perspective process because lab examination problems are operational, academic, and technical at the same time \cite{elicit1}. No single elicitation technique could capture the full picture. The project therefore used interviews, prototyping, and brainstorming sessions to define realistic requirements.

\subsubsection{Interviews}
Interviews are a direct requirement elicitation technique in which stakeholders are asked structured questions to explain their current workflow, difficulties, and expectations. In this project, interviews were used to understand how teachers and academic staff currently schedule practical exams, supervise students, collect files, and review incidents.

Two key questions guided this technique. The first question asked what difficulties arise during practical lab exams when many students must be supervised at the same time. The second question asked what information teachers and administrators need after the exam for attendance verification, result handling, and academic review.

The outcomes of the interviews were clear and actionable. First, they highlighted the need for role-based access and controlled scheduling. Second, they showed that live monitoring and incident logging are essential for exam oversight. Third, they confirmed that teachers need structured reporting after the session rather than scattered manual records.

\subsubsection{Prototyping}
Prototyping is a requirement elicitation technique in which early interface sketches or mock workflows are shown to stakeholders so that hidden usability issues can be discovered before full implementation. In this project, simple interface ideas for login, scheduling, monitoring, and submission were considered to refine the system flow.

Two key questions guided this technique. The first question asked whether the proposed dashboards and screens were understandable enough for use during time-sensitive exam conditions. The second question asked whether the sequence from login to submission matched the actual working pattern of a lab exam.

The outcomes prototyping improved the requirement quality in three ways. First, it clarified the need for a simpler teacher dashboard with visible student status. Second, it confirmed that students need a guided exam flow with clear timing and submission prompts. Third, it exposed the importance of integrating attendance, monitoring, and reporting instead of treating them as unrelated screens.

\subsubsection{Brainstorming Session 1: Academic Workflow Discussion}
Brainstorming is a collaborative elicitation technique in which team members discuss ideas, problems, and possible solutions in multiple sessions to identify realistic and prioritized system requirements. In the first session, the discussion focused on the complete academic workflow, starting from exam creation and ending with final result review.

Two key questions guided this session. The first question asked which steps of the current practical exam workflow consume the most manual effort. The second question asked which modules are essential if the system is to provide value from the very first prototype.

The outcomes of the first brainstorming session established three priorities. First, exam scheduling and role management were identified as foundation modules. Second, lab allocation and attendance were recognized as necessary controls for a real computer-lab environment. Third, the team agreed that result management must remain connected to the earlier monitoring and submission steps.

\subsubsection{Brainstorming Session 2: Monitoring and Reporting Discussion}
The second brainstorming session focused specifically on monitoring, identity verification, incident handling, and reporting. This session explored how AI-based alerts, face recognition, and evidence logging could be made useful within the practical limitations of an FYP project.

Two key questions guided this session. The first question asked which monitoring events should be recorded as meaningful evidence during a live exam. The second question asked what type of final reports would help teachers and administrators review the session efficiently.

The outcomes of the second brainstorming session refined the advanced modules of the project. First, the team prioritized face recognition before exam access rather than after the session. Second, it defined incident logs and alerts as central monitoring outputs. Third, it confirmed that attendance, incidents, and results should appear in one reporting layer so that post-exam review remains coherent.

\section{Time Frame}

The project time frame translates the Agile development plan into a structured schedule consisting of eight sprints. Each sprint lasts four weeks, allowing iterative development, review, and refinement while keeping the project aligned with academic deadlines.

This sprint-based schedule ensures continuous progress, regular supervisor feedback, and balanced development across all twelve modules of the proposed system. The final sprint focuses on integration and polishing so that the prototype is complete and academically presentable.

\begin{table}[H]
\centering
\caption{Project Time Frame (8 Sprints)}
\label{tab:timeframe}
\renewcommand{\arraystretch}{1.25}
\begin{tabular}{lp{8.2cm}c}
\hline\hline
\textbf{Sprint} & \textbf{Modules and Key Activities} & \textbf{Duration} \\
\hline\hline
Sprint 1   & Requirement analysis, backlog preparation, and system planning. & 4 Weeks \\
Sprint 2   & Authentication and Device Binding (Module 1), User and Role Management (Module 7). & 4 Weeks \\
Sprint 3   & Exam Creation and Configuration (Module 8), Eligibility and Access Control (Module 9). & 4 Weeks \\
Sprint 4   & Pre-Exam System Verification and Mock Test (Module 2), Exam Dashboard (Module 3), Secure Exam Environment Module (Module 4). & 4 Weeks \\
Sprint 5   & Live Monitoring Client Service (Module 5), Live Proctoring Dashboard (Module 10). & 4 Weeks \\
Sprint 6   & Exam Submission and Backup (Module 6), AI-Assisted Evaluation Module (Module 11). & 4 Weeks \\
Sprint 7   & Logs, Analytics and Audit Trail System (Module 12). & 4 Weeks \\
Sprint 8   & System integration, end-to-end testing, bug fixing, and final document polishing. & 4 Weeks \\
\hline\hline
\end{tabular}
\end{table}
